%% ==========================================================================
%%  The model and the conjecture --- general binary valuations.
%%
%%  Framing only: the class, the conjecture, what carries over from the two
%%  theorems it generalizes, and the lower bound.  No approach, no algorithm,
%%  no analysis of where a proof would run into trouble.
%%
%%  SELF-CONTAINED.  This document does not \input report/sections/, so every
%%  definition it needs is restated here and every external result is cited
%%  rather than \ref'd.
%%
%%  REFERENCING: theorem-like environments share one counter in ../preamble.tex,
%%  so \Cref renders all of them as "Theorem".  Use explicit
%%  "Lemma~\ref{...}", "Observation~\ref{...}", etc.
%% ==========================================================================
\section{The Model and the Conjecture}
\label{sec:gb-framing}

% --------------------------------------------------------------------------
\subsection{General binary valuations}
\label{sec:gb-model}

Let $\agents = \set{1,\dots,n}$ be a set of agents and $\items$ a finite set of
$m$ indivisible items. Each agent $i \in \agents$ has a valuation
$v_i : 2^{\items} \to \R$ with $v_i(\emptyset) = 0$.

\begin{definition}[General binary valuation]
\label{def:gb-class}
A valuation $v_i$ is \emph{general binary} if every marginal lies in
$\set{-1,0,1}$:
\[
  v_i(S \cup \set g) - v_i(S) \;\in\; \set{-1,0,1}
  \qquad\text{for every } S \subseteq \items \text{ and every } g \in \items \setminus S .
\]
\end{definition}

Nothing further is assumed: $v_i$ need not be additive, submodular,
subadditive, or monotone, and no relation is assumed between the valuations of
different agents. Two consequences of the definition are worth stating
explicitly, since both are excluded in the special cases this class
generalizes.

\begin{itemize}
  \item A single item may be a \emph{good} for one agent and a \emph{chore} for
    another: $v_i(g \mid S) = 1$ while $v_j(g \mid T) = -1$.
  \item For a \emph{single} agent, an item's marginal may be positive on one
    set and negative on another: $v_i(g \mid S) = 1$ and $v_i(g \mid T) = -1$
    for $S \ne T$. Such an item is neither a good nor a chore for that agent
    (Lemma~\ref{lem:gb-not-doubly-monotone}).
\end{itemize}

We write $v_i(g \mid S) := v_i(S \cup \set g) - v_i(S)$ for a marginal.

An \emph{allocation} $\alloc = (\bundle1,\dots,\bundle n)$ is a tuple of
pairwise disjoint subsets of $\items$, assigned one per agent; it is
\emph{complete} if $\bigcup_i \bundle i = \items$. A subsidy vector is
$\subsidy \in \R^n_{+}$, and agents have quasi-linear utilities.

\begin{definition}[Envy-free solution]
\label{def:gb-efsolution}
An allocation $\alloc$ and a subsidy vector $\subsidy$ constitute an
\emph{envy-free solution} $(\alloc,\subsidy)$ if
\[
  v_i(\bundle i) + \subsidy_i \;\ge\; v_i(\bundle j) + \subsidy_j
  \qquad\text{for all } i,j \in \agents .
\]
An allocation is \emph{envy-freeable} if some $\subsidy \in \R^n_+$ makes
$(\alloc,\subsidy)$ an envy-free solution.
\end{definition}

For an allocation $\alloc$, the \emph{envy graph} $\envygraph{\alloc}$ is the
complete weighted digraph on $\agents$ in which the arc $(i,j)$ carries weight
\begin{equation}
\label{eq:gb-arc}
  \edgew{\alloc}(i,j) \;:=\; v_i(\bundle j) - v_i(\bundle i),
\end{equation}
and $\pathw{\alloc}(i)$ denotes the maximum weight of a directed path in
$\envygraph{\alloc}$ starting at $i$, the empty path of weight $0$ included.

% --------------------------------------------------------------------------
\subsection{The conjecture}
\label{sec:gb-conjecture}

\begin{conjecture}
\label{conj:gb-main}
For every instance $\langle \agents, \items, \set{v_i}_{i \in \agents}\rangle$
in which every $v_i$ is general binary, there exist a complete allocation
$\alloc$ and a subsidy vector $\subsidy \in \set{0,1}^n$, with
$\sum_{i \in \agents} \subsidy_i \le n-1$, such that $(\alloc,\subsidy)$ is an
envy-free solution. Moreover $\alloc$ and $\subsidy$ can be computed in
polynomial time given value-oracle access to the $v_i$.
\end{conjecture}

Conjecture~\ref{conj:gb-main} is the common generalization of two theorems.

\begin{observation}[The two boundary cases are known]
\label{obs:gb-boundary}
Restricting Definition~\ref{def:gb-class} to marginals in $\set{0,1}$ gives the
\emph{dichotomous} class, for which Conjecture~\ref{conj:gb-main} is Theorem~4
of Barman, Krishna, Narahari and Sadhukhan~\cite{BKNS22}. Restricting it to
marginals in $\set{-1,0}$ gives the \emph{negative dichotomous} class, for
which Conjecture~\ref{conj:gb-main} is the main theorem of the companion
document \texttt{report/main.tex}, proved for every $n$.
\end{observation}

Both boundary results are stated at the same generality as
Definition~\ref{def:gb-class} within their own sign restriction --- neither
assumes additivity, submodularity or subadditivity --- so
Conjecture~\ref{conj:gb-main} adds exactly one thing to what is known: the two
signs may occur in the same instance.

% --------------------------------------------------------------------------
\subsection{What transfers}
\label{sec:gb-transfers}

The following facts hold for general binary valuations and may be used without
re-derivation.

\begin{observation}[Integrality]
\label{obs:gb-integrality}
Let $v_i$ be general binary. Then $v_i(S) \in \Z$ and $\abs{v_i(S)} \le \abs S$
for every $S \subseteq \items$; every arc weight $\edgew{\alloc}(i,j)$ is an
integer; and for every envy-freeable allocation $\alloc$ the minimal subsidy
vector satisfies $\optsubsidy \in \Z^n_{\ge 0}$.
\end{observation}

\begin{proof}
Fix $i$ and enumerate $S = \set{g_1,\dots,g_k}$. Telescoping from
$v_i(\emptyset) = 0$ writes
$v_i(S) = \sum_{u=1}^{k} v_i\bigl(g_u \mid \set{g_1,\dots,g_{u-1}}\bigr)$, a
sum of $k$ terms each in $\set{-1,0,1}$ by Definition~\ref{def:gb-class};
hence $v_i(S) \in \Z$ and $\abs{v_i(S)} \le k$. Each arc weight
\eqref{eq:gb-arc} is a difference of integers, so the weight of any directed
path is an integer. By Theorem~\ref{thm:gb-hs-minsubsidy},
$\optsubsidy_i = \pathw{\alloc}(i)$ is a maximum of such weights together with
$0$, hence a non-negative integer.
\end{proof}

Observation~\ref{obs:gb-integrality} is what makes
Conjecture~\ref{conj:gb-main} a statement about a \emph{bound} rather than
about rounding: once $\optsubsidy$ is known to be integral, proving
$\optsubsidy_i \le 1$ for every $i$ is the same as proving
$\optsubsidy \in \set{0,1}^n$.

The two structural results of Halpern and Shah~\cite{HS19} transfer verbatim.
Their proofs use only the arc weights \eqref{eq:gb-arc} and permutations of the
bundles; no additivity, no monotonicity and no sign condition on $v_i$ enters,
which is also the generality at which~\cite{BKNS22} invokes them.

\begin{theorem}[\cite{HS19}]
\label{thm:gb-hs-characterisation}
For any allocation $\alloc = (\bundle1,\dots,\bundle n)$ the following are
equivalent.
\begin{enumerate}
  \item[(i)] $\alloc$ is envy-freeable.
  \item[(ii)] $\alloc$ maximises utilitarian welfare across all reassignments
    of its own bundles: $\sum_{i} v_i(\bundle i) \ge \sum_{i} v_i(A_{\sigma(i)})$
    for every permutation $\sigma$ of $\agents$.
  \item[(iii)] $\envygraph{\alloc}$ has no positive-weight directed cycle.
\end{enumerate}
\end{theorem}

\begin{theorem}[\cite{HS19}]
\label{thm:gb-hs-minsubsidy}
For any envy-freeable allocation $\alloc$, the subsidy vector $\optsubsidy$
given by $\optsubsidy_i := \pathw{\alloc}(i)$ realises an envy-free solution
$(\alloc,\optsubsidy)$, and $\optsubsidy_i \le \subsidy_i$ for every $i$ and
every $\subsidy$ with $(\alloc,\subsidy)$ an envy-free solution.
\end{theorem}

\begin{proposition}[Some agent is unpaid]
\label{prop:gb-zero-coordinate}
For every envy-freeable allocation $\alloc$ there is an agent $h$ with
$\optsubsidy_h = 0$. Consequently, if $\optsubsidy \in \set{0,1}^n$ then
$\sum_{i \in \agents} \optsubsidy_i \le n-1$.
\end{proposition}

\begin{proof}
By Theorem~\ref{thm:gb-hs-characterisation}(iii) the graph
$\envygraph{\alloc}$ has no positive-weight cycle, so a walk of maximum weight
may be taken simple and $\pathw{\alloc}$ is finite; it is non-negative because
the empty path qualifies. Choose a path $P$ attaining
$\max_{i} \pathw{\alloc}(i)$ and let $h$ be its final vertex. If
$\pathw{\alloc}(h) > 0$ there is a positive-weight walk $Q$ leaving $h$, and
$P$ followed by $Q$ is a walk from $P$'s start of weight
$\edgew{\alloc}(P) + \edgew{\alloc}(Q) > \edgew{\alloc}(P)$, contradicting the
maximality of $P$ among walks. Hence $\pathw{\alloc}(h) = 0$. The second claim
follows because at most $n-1$ coordinates of $\optsubsidy$ can then equal $1$.
\end{proof}

Proposition~\ref{prop:gb-zero-coordinate} means the bound on the total in
Conjecture~\ref{conj:gb-main} is not an independent requirement but a
consequence of the per-agent one, so the conjecture is the single-clause
statement
\[
  \exists\, \alloc \ :\quad \max_{i \in \agents} \pathw{\alloc}(i) \;\le\; 1 .
\]

% --------------------------------------------------------------------------
\subsection{Tightness}
\label{sec:gb-tightness}

\begin{example}[The $n-1$ lower bound]
\label{ex:gb-lowerbound}
Take $n$ agents and a single item $g$ with $v_i(\set g) = 1$ for every $i$.
Every $v_i$ is general binary. Under any complete allocation one agent holds
$g$ and the other $n-1$ hold nothing; each of those $n-1$ agents envies the
holder by exactly $1$, so every envy-free solution assigns each of them a
subsidy at least $1$ more than the holder's. The minimum total is therefore
$n-1$.
\end{example}

The same bound is forced from the opposite sign: $n$ agents and $m = n-1$ items
with $v_i(S) = -\abs S$ for every $i$ leaves some agent empty-handed under any
complete allocation, and each of the $n-1$ loaded agents envies that agent by
$1$. Either way, if Conjecture~\ref{conj:gb-main} holds then its bound is
tight.

% --------------------------------------------------------------------------
\subsection{General binary is not a doubly monotone class}
\label{sec:gb-not-doubly-monotone}

Recall that an instance is \emph{doubly monotone} when, for each agent $i$ and
each item $g$, either $v_i(g \mid S) \ge 0$ for every $S$ (the item is a
\emph{good} for $i$) or $v_i(g \mid S) \le 0$ for every $S$ (a \emph{chore}
for $i$). Both boundary classes of Observation~\ref{obs:gb-boundary} are
doubly monotone; their common generalization is not.

\begin{lemma}
\label{lem:gb-not-doubly-monotone}
There is a general binary valuation that is not doubly monotone.
\end{lemma}

\begin{proof}
Let $\items = \set{a,b,g}$ and define
\[
  v(S) := \begin{cases} \abs S & \text{if } \abs S \le 2,\\ 1 & \text{if } S = \items. \end{cases}
\]
Then $v(\emptyset) = 0$, and every marginal lies in $\set{-1,0,1}$: adding an
item to a set of size $0$ or $1$ raises the value by exactly $1$, and adding
the third item to a set of size $2$ lowers it from $2$ to $1$. So $v$ is
general binary. But $v(g \mid \emptyset) = 1 > 0$ while
$v(g \mid \set{a,b}) = 1 - 2 = -1 < 0$, so $g$ is neither a good nor a chore
for this agent.
\end{proof}

\begin{remark}
\label{rem:gb-no-baseline}
Lemma~\ref{lem:gb-not-doubly-monotone} has a concrete consequence for what may
be assumed. The reduction of Kawase, Makino, Sumita, Tamura and
Yokoo~\cite{KMSTY24}, which converts any EF1 allocation into an envy-free one
and yields $n-1$ per agent, is stated for doubly monotone valuations under a
two-sided normalisation. It therefore applies to each boundary class of
Observation~\ref{obs:gb-boundary} separately --- and did supply the baseline
against which the negative dichotomous result was measured --- but it does
\emph{not} apply to general binary instances off the shelf.
\end{remark}
